\chapter{اللقاء الثامن والعشرون: الخمر والميسر والإنفاق وأحكام اليتامى}

\section{مقدمة}
السلام عليكم ورحمة الله وبركاته. الحمد لله رب العالمين، وأشهد أن لا إله إلا الله وحده لا شريك له، وأشهد أن محمداً عبده ورسوله. اللهم صل على محمد وعلى آل محمد كما صليت على آل إبراهيم إنك حميد مجيد، اللهم بارك على محمد وعلى آل محمد كما باركت على آل إبراهيم إنك حميد مجيد.

ختم الشيخ في هذا اللقاء المجلد الثالث من تفسير \rawi{الطبري}، ووقف عند جملة من الآيات الجامعة في سورة البقرة: آية الخمر والميسر، وآية الإنفاق، وآيات اليتامى، مع عناية خاصة ببيان تدرج التحريم، وسعة مفهوم الميسر، ودقة \rawi{الطبري} في جمع الأقوال ورد بعضها إلى بعض، وبيان أثر أسماء الله الحسنى في ختام الآيات.

\separator

\section{تأويل (يسألونك عن الخمر والميسر)}
\isnad{قال أبو جعفر:} الخمر كل شراب خامر العقل فغطاه وستره، والميسر هو القمار وما كان في معناه من المغالبة على المال أو العوض بغير حق.

ثم ساق \rawi{الطبري} آثار السلف في تفسير الميسر، فظهر من مجموعها سعة دلالته عندهم:
\begin{itemize}
    \item فالقمار الصريح داخل فيه.
    \item وما كان على المخاطرة في الأموال داخل فيه.
    \item وحتى بعض ألعاب الصبيان إذا قامت على معنى المغالبة والمخاطرة ألحقت به.
\end{itemize}

ونبّه الشيخ إلى أن هذا الموضع من أوائل ما عيبت به الخمر في القرآن قبل نزول التحريم القاطع في سورة المائدة، فهو من مواضع التدرج التشريعي الواضحة.

وشدد هنا على أن إضافة الإثم إلى الخمر والميسر في هذه الآية لا تعني أن التحريم النهائي قد وقع بها، بل المراد بيان ما ينشأ عنهما من آثام ومفاسد تمهد النفوس لتركهما، حتى جاء النص القاطع بعد ذلك في سورة المائدة.

\begin{fawaid}[التحريم قد يسبقه تمهيد يربي النفوس على ترك المحرم]
من حكمة الشريعة أن بعض المحرمات العظيمة لم ينزل فيها التحريم دفعة واحدة، بل سبقته مراحل تبين قبحها وفسادها، حتى تتهيأ النفوس لقبول الحكم الكامل، كما في الخمر.
\end{fawaid}

\section{تأويل (قل فيهما إثم كبير ومنافع للناس)}
ذكر \rawi{الطبري} أن في الخمر والميسر منافع دنيوية معروفة عند الناس:
\begin{itemize}
    \item فمنفعة الخمر فيما يجدونه من اللذة، وفي أثمانها قبل تحريمها.
    \item ومنفعة الميسر فيما يصيبونه من أنصباء الجزور أو ما يربحه بعضهم في صورة الظفر.
\end{itemize}

لكنه رجح أن الإثم الكبير فيهما أعظم من هذه المنافع، وفسر أعظم إثم الخمر بزوال العقل حتى يعزب العبد عن معرفة ربه وذكره، وفسر أعظم ما في الميسر بما يجره من الظلم، وأكل المال بالباطل، والصد عن ذكر الله والصلاة، وإيقاع العداوة والبغضاء.

وفي هذا الموضع بيّن الشيخ أن وجود بعض النفع في الشيء لا يمنع تحريمه إذا كانت مفسدته أغلب وأعظم.

ونبّه أيضاً إلى فائدة أصولية مهمة: أن هذه الآية لا يسوغ أن تختزل في عبارة مجملة من نحو: ``وجود المفسدة يكفي في المنع''؛ بل الذي دلت عليه هو الموازنة بين قدر المصلحة والمفسدة، وأن الحكم يتبع غلبة أحد الجانبين وعظمه.

\begin{fawaid}[وجود المنفعة لا يقتضي الإباحة إذا غلبت المفسدة]
ليس كل ما فيه نفع يباح للناس، بل الميزان الشرعي أن تنظر المصلحة والمفسدة جميعاً. فإذا كانت المفسدة أعظم وأعم، لم يرفع وجود بعض المنافع حكم التحريم.
\end{fawaid}

\separator

\section{تأويل (ويسألونك ماذا ينفقون قل العفو)}
انتقل \rawi{الطبري} إلى تفسير \textit{العفو} في هذه الآية، وبيّن أنه ما فضل عن الحاجة والكفاية مما لا يجحف بصاحبه إذا أنفقه.

وفي هذا الموضع تتجلى طريقة القرآن في ضبط باب الإنفاق: فلا إسراف يضر المنفق ومن يعول، ولا شح يمنع الحق الواجب والمندوب.

\section{تأويل (كذلك يبين الله لكم الآيات لعلكم تتفكرون في الدنيا والآخرة)}
فسر \rawi{الطبري} الآية بأن الله يبين لعباده أحكامه ودلائل هدايته ليوقع لهم التفكر الصحيح في أمر الدنيا والآخرة، فلا ينظروا إلى العاجل مجرداً عن العاقبة، ولا إلى المصالح الجزئية دون مآلاتها.

\separator

\section{تأويل (ويسألونك عن اليتامى قل إصلاح لهم خير)}
\isnad{قال أبو جعفر:} لما نزل التشديد في أكل أموال اليتامى، شق على الناس أمر مخالطتهم، فسألوا النبي صلى الله عليه وسلم عن ذلك، فأنزل الله التخفيف: \begin{ayah}قُلْ إِصْلَاحٌ لَهُمْ خَيْرٌ\end{ayah}.

فدل ذلك على أن المقصود الأعظم في باب اليتيم ليس مجرد العزل الشكلي لماله عن مال وليه، بل تحقيق الإصلاح له وحفظ ماله وتنميته.

وقد جمع الشيخ هنا بين الروايات الواردة في سبب السؤال، وردها إلى أصل واحد: أن المسلمين لما خافوا من الوعيد في أموال اليتامى اشتد عليهم العزل والمجانبة، سواء أكان ذلك بعد سؤال صريح، أم بمشقة نزل القرآن برفعها، أم على وفق عادة كانت عند العرب في التحرج من مال اليتيم؛ فأنزل الله الحكم الذي يجمع بين الحفظ والتيسير.

\section{تأويل (وإن تخالطوهم فإخوانكم)}
بيّن \rawi{الطبري} أن الله أباح مخالطة اليتامى في المطعم والمشرب والمسكن ونحو ذلك إذا كان القصد الإصلاح لا الإفساد. وذكر الشيخ هنا فرقاً لطيفاً في الإعراب والمعنى بين قوله: \textit{وإن تخالطوهم فإخوانكم}، وبين مواضع أخرى يجيء فيها النصب على الحال أو على معنى الأمر؛ لأن المعنى هنا: فهم إخوانكم على كل حال، خالطتموهم أو لم تخالطوهم، فلا يراد إنشاء أخوة جديدة بالمخالطة، بل تقرير أخوة قائمة تقتضي الرفق والإعانة.

ثم أكد أن الآية لا تفتح باب التحايل على أموال اليتامى باسم المخالطة أو الولاية، بل تقرر أصل الأخوة والرفق مع بقاء المراقبة الشديدة للنية والقصد.

\begin{fawaid}[الله يعلم المفسد من المصلح ولو تشابهت الصور الظاهرة]
قد يكون العمل في ظاهره مشروعاً أو مباحاً، لكن القصد الذي وراءه يفسده ويجعله وبالاً على صاحبه. ولذلك كانت هذه الآية من أعظم ما يردع من يتخذ الأعمال المشروعة أو الدعوية أو المالية ستاراً لمقاصد فاسدة.
\end{fawaid}

\section{تأويل (والله يعلم المفسد من المصلح)}
وقف الشيخ عند هذه الآية وقفة تربوية عظيمة، وقرر أن من أخطر ما يبتلى به الناس أن يدخلوا أبواباً مشروعة أو مباحة ونياتهم فيها فاسدة، يطلبون بها أكلاً للمال، أو سمعة، أو تسلطاً، أو إفساداً. وهذه الآية تقطع عليهم هذا الباب:
\begin{itemize}
    \item فالله يعلم حقيقة القصد وإن خفي على الناس.
    \item والعمل الذي بني على نية فاسدة لا يكون صالحاً وإن تشبه بصورة الصلاح.
    \item والله لا يصلح عمل المفسدين وإن بدا لهم في أوله أنه نافع.
\end{itemize}

\begin{fawaid}[فساد النية يفسد العمل وإن كان في صورته حسناً]
ليست العبرة بمجرد صورة العمل، بل بحقيقته الباطنة ومقصوده. فمن دخل باباً مشروعاً وقصده فيه الإفساد أو الاستطالة أو أكل الحقوق، لم تنفعه صورة العمل، لأن الله يعلم المفسد من المصلح.
\end{fawaid}

\separator

\section{تأويل (ولو شاء الله لأعنتكم)}
ذكر \rawi{الطبري} أن معنى \textit{لأعنتكم}: لأشق عليكم، أو لأحرجكم، أو لضيق عليكم، أو لأوقعكم في المشقة والعنت لو حرم عليكم هذه المخالطة وضيّق عليكم في شأن اليتامى.

ثم نبه الشيخ إلى أن هذا الموضع من أحسن أمثلة اختلاف التنوع؛ إذ اختلفت ألفاظ السلف في تفسير \textit{العنت}، لكنها متقاربة المعنى، ترجع كلها إلى الشدة والحرج والمشقة.

وزاد أن الآية تبرز كمال الرحمة في التشريع؛ فالله قادر على أن يضيق ويشدد، لكنه لم يشأ ذلك، بل فتح باب المخالطة بضابط الإصلاح رحمة بعباده، ليقوم حق اليتيم من غير أن ينقلب التكليف إلى حرج يعسر احتماله.

\begin{fawaid}[اختلاف الألفاظ لا يعني دائماً اختلاف المعاني]
قد ينقل المفسر عن السلف عبارات متعددة في تفسير اللفظ الواحد، كقول بعضهم: لأشق عليكم، أو لأحرجكم، أو لأضيق عليكم. وهذه ليست أقوالاً متباينة على الحقيقة، بل هي من اختلاف التنوع الذي يجمعه معنى كلي واحد.
\end{fawaid}

\section{تأويل (إن الله عزيز حكيم)}
ختم \rawi{الطبري} الآية ببيان مناسب لمعنى الاسمين هنا:
\begin{itemize}
    \item فهو \textit{عزيز} لا يمنعه أحد لو شاء أن يشق على عباده أو يعاقبهم.
    \item وهو \textit{حكيم} لا يشرع إلا ما تقتضيه الحكمة والرحمة والمصلحة.
\end{itemize}

ونبّه الشيخ إلى أن \rawi{الطبري} لا يفسر أسماء الله في خواتيم الآيات تفسيراً معزولاً عاماً فقط، بل يربط كل اسم بسياق الآية التي ختم بها، ويبين أثره في معنى الحكم المذكور.

\begin{fawaid}[أسماء الله في ختام الآيات مرتبطة بسياقها]
من دقة التفسير أن ينظر في مناسبة الاسم الإلهي لخاتمة الآية، لا أن يفسر الاسم مجرداً عن السياق. والطبري من أئمة هذا الباب؛ يبين لماذا ختمت الآية بـ \textit{عزيز حكيم} أو غيرهما، وما أثر ذلك في فهم الحكم.
\end{fawaid}

\separator

ختم الشيخ اللقاء بالتأكيد على أن هذه الآيات تجمع بين الفقه والتربية والمنهج: فهي تبين تدرج التشريع، وضابط الموازنة بين المصالح والمفاسد، وسعة علم الله بالنيات، ورحمته في التخفيف عن عباده، ودقة السلف في اختلاف التنوع، وعظمة القرآن في ربط الأحكام بأسماء الله وصفاته.
